\documentclass[11pt,a4paper]{article}
\usepackage[latin1]{inputenc}
\usepackage{amsmath}
\usepackage{amsfonts}
\usepackage{amssymb}
\usepackage{mathrsfs}
\author{Miguel Angel Rodriguez Feregrino}
\title{Modelo Babesia bigemina}
\begin{document}

\maketitle

Primer modelo 

\begin{equation} 
\left\{\begin{array}{l}
\dfrac{d\mathscr{B}}{dt}
=
\alpha(t)- \beta \mathscr{B}\mathscr{E} - \omega \mathscr{B} + \mu [\rho(1+\phi ) \mathscr{I}]
\\
\dfrac{d\mathscr{E}}{dt}
=
\gamma (t) - \beta \mathscr{B}\mathscr{E} -\rho \mathscr{E} 
\\
\dfrac{d\mathscr{I}}{dt}
=
\beta \mathscr{B}\mathscr{E} - \rho (1+\phi )\mathscr{I}
\\
    \end{array}\right.
\end{equation}

Donde se describe la dinamica entre las poblaciones de Babesias infecciosas libres ($\mathscr{B}$), eritrocitos sanos en circulacion ($\mathscr{E}$) y eritrocitos infectados ($\mathscr{I}$).

Las reacciones por tanto se pueden describir como:

\begin{itemize}
\item \textbf{Reaccion 1}: Introduccion de babesias infectivas al sistema. $\alpha(t)$
\item \textbf{Reaccion 2}: Introduccion de eritrocitos sanos al sistema. $\gamma (t)$
\item \textbf{Reaccion 3}: Infeccion de eritrocitos sanos, accion de la interaccion con babesias infectivas. $\beta \mathscr{B}\mathscr{E}$
\item \textbf{Reaccion 4}: Muerte natural de babesias infectivas en circulacion. $\omega \mathscr{B}$
\item \textbf{Reaccion 5}: Muerte natural de eritrocitos sanos. $\rho \mathscr{E}$
\item \textbf{Reaccion 6}: Muerte de eritrocitos infectados, dado por la muerte natual de los eritrocitos mas la accion de la infeccion en esta. $\rho (1+\phi )\mathscr{I}$
\item \textbf{Reaccion 7}: Liberacion de babesias infectivas al ambiente resultado de la muerte de eritrocitos infectados. $\mu [\rho(1+\phi ) \mathscr{I}]$

\end{itemize}

 
\end{document}

